\section{Proofs for the main text}\label{app:proofs}

This appendix contains the proofs omitted from the main text, in the order of the sections they belong to.
Short arguments given in full in the main text are not repeated.

\subsection{Owned choice and rational play}\label{app:proofs-rational}

\begin{proof}[Proof of \cref{game:thm:lower}]
Zero occurrences cannot express dependence on that input. Suppose $x$
occurs once. Read the unique path from its occurrence to the root, using
commutativity to put the current value first. It has the form
\[
 X\longmapsto X\ow{r_1}g_1(Y)
   \longmapsto\cdots\longmapsto (\cdots)\ow{r_n}g_n(Y),
\]
where each $g_j(Y)$ is nonempty and independent of $X$. We allow arbitrary
such functions, which only strengthens the obstruction.

For fixed $Y$, if some $g_j(Y)=\{r_j\}$, that step erases all dependence
on $X$, and later steps cannot recover it. For $Y=\{r\}$ this is
incompatible with $X\cup\{r\}$, which varies with $X$. Therefore no such
side value equals its owner singleton. Every step then preserves all
members of its current input: if the input is the owner singleton it
remains so, and otherwise the operation takes a union.

If the path is empty, its value is $X$, not union. Otherwise let $r$ be
its innermost owner and set $Y=\{r\}$. Its first side value differs from
$\{r\}$ and is nonempty, so contains some $k\ne r$. Choose
$j\notin\{r,k\}$. Starting with $X=\{j\}$, the first step introduces
$k$. No later step can remove it. The required result $\{j,r\}$ excludes
$k$, a contradiction. Interchange the variables for the other bound.
\end{proof}

The syntax-to-path correspondence is the direct decomposition of a term at its unique input occurrence, as used in
this proof.

\subsection{Player views}\label{app:proofs-views}

\begin{proof}[Proof of \cref{view:thm:view}]
For nonempty $S,T$, direct calculation gives
\[
 S\ow r T=\{i\}\ \Longleftrightarrow\
 \begin{cases}
 S=\{i\}\text{ or }T=\{i\},&r=i,\\
 S=T=\{i\},&r\ne i.
 \end{cases}
\]
Also, membership of the owner $r$ is simply the disjunction of its
membership in the two inputs. For $k\ne r$, membership of $k$ is
\[
 (k\in S\text{ or }k\in T)\ \land\ S\ne\{r\}\ \land\ T\ne\{r\}.
\]
These equations show that the right-hand side of~\eqref{view:eq:view} is a
congruence, hence is preserved by every context.

Conversely, the identity context tests membership of $k$. For $i\ne k$,
the context $x\ow i\{k\}$ excludes $k$ exactly when $x=\{i\}$. These
contexts distinguish all the listed cases. The class $\nv$ is inhabited
because there are two distinct roles other than $k$.
\end{proof}

Ground constants suffice to represent arbitrary polynomial parameters because union realizes every nonempty subset.

\begin{proof}[Proof of \cref{view:thm:tests}]
\emph{Sufficiency.} Suppose $f,g$ differ at a valuation $v$, and choose
$k$ belonging to exactly one output. Choose an ambiguous $F\in D$
omitting $k$. Replace each input set $S$ as follows:
\[
 S\longmapsto
 \begin{cases}
 \{k\},&k\in S,\\
 \{i\},&S=\{i\},\ i\ne k,\\
 F,&\text{otherwise}.
 \end{cases}
\]
The replacement belongs to $D$ and has the same $k$-view. All operations
preserve that view, so the outputs still differ at $k$. The replacement
may use the \emph{same} $F$ for every variable. This proof works at any
arity, not just for unary functions.

\emph{Necessity of singleton tests.} The unary classification in \cref{game:thm:unary} realizes all functions
\[
 \Reg(C,E)(S)=\begin{cases}\{i\}\cup E_i,&S=\{i\},\\
 S\cup C,&|S|\geq2,\end{cases}
 \quad E_i\subseteq C\setminus\{i\}.
\]
Fix $i$ and $k\ne i$. Take $C=I$, all rows full except
$E_i=I\setminus\{i,k\}$. This function differs from the constant $I$
only at $\{i\}$. Hence that singleton must be tested.

\emph{Necessity of ambiguity coverage.} For each $k$, the same
classification realizes
\[
 F_k(S)=\begin{cases}S,&|S|=1,\\S\cup\{k\},&|S|\geq2.\end{cases}
\]
This differs from the identity exactly on ambiguous sets omitting $k$.
At least one such input must be tested.
\end{proof}

\begin{proof}[Proof of \cref{view:cor:optimal}]
Every singleton is necessary. One ambiguous set has nonempty
intersection with itself, so at least two are needed. Two disjoint pairs
suffice for $K\geq4$. At $K=3$, every ambiguous set has size at least
two, so any two intersect. All three pairs are necessary and sufficient.
\end{proof}

\subsection{Pair tests}\label{app:proofs-tests}

\begin{proof}[Proof of \cref{sym:lem:provenance}]
A winner or singleton-status input contributes a constant. An ambiguous
input contributes its identity permutation. A role switch transports
$C$ and composes the permutations. A nonempty $P_j$ contributes an
ambiguous set, so that expression cannot equal an owner singleton. Thus
an owned choice collapses exactly when some child is a constant owner
singleton on the cell. Otherwise union its constants and coefficients.
This proves the representation by induction. It is not asserted to be
canonical: distinct coefficients may have the same action on ambiguous
inputs.
\end{proof}

The representation in \cref{sym:lem:provenance} is not asserted to be canonical: distinct coefficients may have the same
action on ambiguous inputs. The necessity of pairs uses the following gadget.

\begin{lemma}[Resetting atomic inputs]\label{sym:lem:reset}
If a unary term function $F$ is extensive, start with $U(x)=F(x)$ and
successively replace $U(x)$ by $U(x)\ow i x$, for each $i\in I$. This gives a term
$\operatorname{Reset}(F)$ which fixes all singletons and agrees with $F$
on every ambiguous input.
\end{lemma}
\begin{proof}
At input $\{i\}$, its own suppression returns $\{i\}$. A different
owner preserves the current value because it contains $i$. On ambiguous
input $S$, the current value contains $S$; neither operand is a
singleton, and their union is still the current value. Extensivity is
preserved at each step.
\end{proof}

\begin{proof}[Proof of \cref{sym:thm:optimal}]
Sufficiency is Theorem~\ref{sym:thm:pairs}. Each singleton is necessary
already in the reduct: its unary synthesis gives a term differing from
the constant $I$ only at that singleton. These separators are now
kernel-checked in \texttt{Synthesis.lean}.

For necessity of a particular pair $P$, choose $k\notin P$, put
$C=I\setminus\{k\}$, and form the extensive term
\[
 F(x)=x\cup C\cup\bigcup_{i\in I\setminus(P\cup\{k\})}(i\ k)(x).
\]
Here $(i\ k)$ is a transposition. If $x\subseteq P$, then $F(x)=C$.
Otherwise $x$ contains $k$, or contains an $i$ whose transposition
supplies $k$, and $F(x)=I$. Thus, on ambiguous inputs, $F$ differs from
$I$ only at $P$. By Lemma~\ref{sym:lem:reset},
$\operatorname{Reset}(F)$ and $\operatorname{Reset}(I)$ agree at every
singleton and every ambiguous input other than $P$, and disagree at
$P$. Every complete input domain must therefore contain $P$.
\end{proof}

\subsection{The intrinsic core}\label{app:proofs-core}

\begin{proof}[Proof of \cref{gen:thm:core}]
On the all-ambiguous cell, track a constant set $C$ and the
variable-occurrence permutations. If any winner leaf occurs, $C$ is
nonempty: it starts nonempty at that leaf, is transported by switches,
survives unions, and a collapse at an owned node replaces it by a
nonempty owner singleton. A node with no constant contribution cannot
collapse on this cell, since every live child is ambiguous. Thus a term
containing a winner leaf has some fixed nonempty $C$ in every ambiguous
diagonal output. But $\bigcap_{|S|\geq2}S=\varnothing$ for $K\geq3$,
so the term cannot be idempotent.

Without winner leaves, every intermediate value on this cell is
ambiguous. Every choice is union. The diagonal value is therefore the
union of $\pi S$ over all variable occurrences. Idempotency forces each
$\pi S\subseteq S$. If $\pi(i)\ne i$, a pair containing $i$ and omitting
$\pi(i)$ contradicts that inclusion. Every accumulated permutation is
identity.

Transport switches through choices using
$\pi(u\star_r v)=\pi u\star_{\pi(r)}\pi v$ and compose them.
All leaf permutations disappear. The remaining constant-free owned term
is idempotent because its primitive choices are. A traversal accumulating
permutations both checks the criterion and transports the owners,
without expanding the tree.
\end{proof}

The formal proof uses persistence of a constant contribution on each status cell and survival of every leaf contribution on
the all-ambiguous cell. The extraction accumulates permutations while traversing the original term and transports owners at
internal nodes.

\begin{proof}[Proof of \cref{gen:cor:trace}]
In the constant-free core all intermediate ambiguous values are
ambiguous; the result is union of the occurring variables. Distinct
supports are separated by a pair and $I$. Composition unions supports,
and every nonempty support is realized by derived union. A nonempty
join cannot satisfy near-unanimity: put the larger of those two values
at a selected argument and the smaller at all others. The Mal'cev
identities would require support both in the first argument alone and
in the third alone. Operations satisfying either collection of
identities are already idempotent. The same two-value test detects the
presence of each argument, proving essentiality. Regularity of identities
is the standard two-element semilattice observation: interpret every
primitive choice as join on $\{P,I\}$, for a proper ambiguous pair $P$.
An identity must then have the same variable support on both sides.
This is the usual regular-identity setting of semilattice sums
\cite{Plonka1967}.
\end{proof}

\begin{proof}[Proof of \cref{gen:cor:views}]
Each $q_k$ is a congruence quotient for every owned choice by \cref{view:thm:view}. Theorem~\ref{gen:thm:core} extends this to all
idempotent operations, even when their original definitions contain
internal switches. All membership bits are recovered from the $\pv$
states, so the family separates values. Each view is onto. Omitting
$q_k$ identifies $I\setminus\{k\}$ and $I$, which every other player
views as $\pv$.
\end{proof}

\begin{proof}[Proof of \cref{gen:thm:idem-domain}]
If $k\in\bigcap D$, the nonidentity term $x\cup\{k\}$ fixes every test.
If two roles have the same incidence vector, their transposition fixes
every test. These necessary conditions apply also to singleton tests.

Conversely suppose every test is ambiguous and the two conditions hold.
An agreeing term's constant contribution lies in $\bigcap D$, so is
empty. Hence it has no winner leaves. Each occurrence permutation
satisfies $\pi S\subseteq S$ on each test, and finiteness makes these
equalities. A permutation fixing each test can only exchange roles with
identical incidence vectors. It is identity. Apply Theorem~\ref{gen:thm:core}.

With $m$ tests the conditions require $K$ distinct binary words of
length $m$, excluding the all-one word. Thus $K\leq2^m-1$. For $K\geq4$
set $m=\lceil\log_2(K+1)\rceil$. Start the role signatures with the $m$
distinct two-bit words supported on a cycle on $m$ bits. Each bit occurs
twice, making every test ambiguous. Add distinct words other than the
all-one word until there are $K$ signatures.

For $K=3$, two tests satisfying the necessary conditions have signatures
$00,01,10$, hence are two distinct singletons. The reduct's ambiguity
adder from \cref{game:thm:unary} fixes every singleton but adds $c$ to
ambiguous inputs, so those tests cannot certify idempotency.
Three pairs certify it.
\end{proof}

\begin{proof}[Proof of \cref{gen:thm:S2}]
Let $\iota_3\subseteq I^3$ be the relation of triples with a repeated
entry. Three input environments having this property in every
coordinate cannot have three distinct singleton outputs:
Theorem~\ref{gen:thm:source} would give one permutation sending an input
triple with repetition to three distinct roles. Every induced function
therefore preserves $\iota_3$.

Use the established presentation $\operatorname{Pol}(\iota_3)=S_2$,
recorded by Szendrei\cite[Proposition~2.2]{Szendrei2012} in the
S\l upecki--Burle classification. It applies for every finite $K\geq3$.
This is a use of that classification, not a new classification of $S_2$.

Conversely every unary map $h:I\to I$ has an atomic realization.
The term $x\star_a(a\ b)(x)$ merges atomic $b$ into $a$, fixing the
others. Merge each $h$-fiber into a representative, then permute the
representatives to the prescribed images.

For distinct output roles $a,b$, tests of input labels are unary maps
into $\{a,b\}$. On these atomic values, $\star_a$ is Boolean OR and
$\star_b$ Boolean AND, with truth $a$. Finite DNF of label tests
realizes every $I^n\to\{a,b\}$. Constants cover the one-valued case.
\end{proof}

\subsection{Reset}\label{app:proofs-reset}

\paragraph{The binary fragment.}
Work first with constant-free, switch-free owned-choice terms, which are automatically idempotent.

\begin{lemma}[One-constant rigidity]\label{reset:lem:rigidity}
A unary owned-choice expression whose only constant is $\{r\}$ is
either constantly $\{r\}$ or inflationary: it contains its input on
every valuation. Consequently, returning $\{r\}$ at any other singleton
$\{q\}$ already forces it to be constant.
\end{lemma}
\begin{proof}
Induct on expressions. Combining two inflationary expressions preserves
inclusion of the input: if an owner wins outright, that singleton
already contains the input. Combining an inflationary expression with
the constant $\{r\}$ returns the constant when the owner is $r$;
otherwise it remains inflationary, because the constant operand cannot
give a different owner an outright win. Two constant expressions remain
constant. An inflationary function cannot send $\{q\}$ to $\{r\}$ when
$q\ne r$. The dichotomy and this rigidity consequence are
kernel-checked in \texttt{OneConstant.lean}.
\end{proof}

\begin{proof}[Proof of \cref{reset:thm:binary}]
The atomic restriction of an idempotent atom-preserving owned term is
a projection. This is \cref{gen:thm:S2}, using $\operatorname{Pol}(\iota_3)=S_2$
\cite[Proposition~2.2]{Szendrei2012}, and an idempotent operation on
$K\geq3$ atoms has image size at least three, so the essentially unary
case of $S_2$ must apply.

Suppose the atomic restriction projects $x$. If $y$ does not occur,
idempotency makes the whole term $x$. Otherwise both variables occur.
Freeze $x=\{r\}$, leaving a unary expression in $y$ with only that
constant. Its atomic outputs are $\{r\}$, so
Lemma~\ref{reset:lem:rigidity} makes it constant on all $y$.

Now let $x$ be ambiguous. If $y$ is ambiguous too, every node computes
union, giving $x\cup y$. If $y=\{r\}$, specialize instead in $x$.
Its atomic restriction is the identity, so rigidity makes it
inflationary. Its result is contained in $x\cup\{r\}$, since owned
choice never introduces a winner absent from its inputs. It contains
$r$ as well: under this valuation every singleton subexpression is
$\{r\}$, by induction, so no different owner can suppress the
contribution of an occurring $y$. The value is therefore
$x\cup\{r\}$. These cases give exactly $y\del x$.
The opposite projection gives the other two operations, all of which
are definable.
\end{proof}

The binary theorem uses the idempotent-elimination theorem and the external atomic-clone classification; its complete proof is
not yet a Lean theorem.

\paragraph{Selection and the chain.}

\begin{proof}[Proof of \cref{reset:lem:selection}]
If the right child is resolved, reset returns it. If it is ambiguous,
reset unions the children and remains ambiguous. This proves the first
three statements by induction. Monotonicity follows from the positive
recursion. For self-independence, at a node with root $j$ that input is
always selected. Otherwise its new guard concerns a different position,
and both child selection predicates are independent of $j$ by induction.
\end{proof}

\begin{proof}[Proof of \cref{reset:thm:three}]
Equal tables give equal selectors and therefore equal values.
If two tables differ, some selector inclusion differs at $B,j$.
Assign $T$ to position $j$, $U$ to the other positions in $B$, and $Z$
elsewhere. Setting $j$ to ambiguous does not change whether $j$ is
selected, by Lemma~\ref{reset:lem:selection}. The marker $c$ occurs in the
result exactly when that position is selected. Thus the terms differ.
Distinct minimal antichains give such $B,j$: take a minimal prerequisite
on one side not containing a prerequisite from the other. All operations
on finite guard tables are polynomial.
\end{proof}

\begin{proof}[Proof of \cref{reset:cor:generator}]
Theorem~\ref{reset:thm:three} identifies their identities. The three-element
algebra is nonassociative:
$(2\del0)\del1=1$, while $2\del(0\del1)=2$.
Every idempotent binary operation on two elements is one of the two
projections or two semilattice operations; all four are associative.
Any generator would have to satisfy the valid identity $x\del x=x$,
so a smaller generator is impossible.
\end{proof}

\paragraph{Path normal forms.}
The same source already develops common-root path normal forms and a
suffix criterion for comparing them
\cite[Lemmas~3.5--3.6]{FreeseEtAl2002}. The correspondence with guards
is exact. For a left-associated path $p$ with pairwise distinct input
positions $i_1,\ldots,i_m$, let
\[
 \sigma_p(i_a)=\{i_{a+1},\ldots,i_m\}.
\]
Its root is $i_m$. Unfolding the selector recursion gives
$i_a\in E_p(B)$ iff $\sigma_p(i_a)\subseteq B$: an occurrence
contributes precisely when every variable to its right is ambiguous.
If $w$ is a product of paths with that common root, its selectors are
the unions of the path selectors. Its canonical guards are therefore
\[
 \mathcal G_j(w)=\min_{\subseteq}
     \{\sigma_p(j):p\text{ is a constituent path of }w,\ j\in p\}.
\]
For another path $p$ with the same root, requiring
$E_p(B)\subseteq E_w(B)$ for every $B$ is equivalent to requiring,
for each occurrence of $j$ in $p$,
some suffix for $j$ in $w$ contained in $\sigma_p(j)$. This is the
suffix-inclusion criterion of their Lemma~3.6, expressed through our
monotone selection predicates. Their order inequality $w\leq p$ reads
$w\del p=w$; for common-root terms this is exactly the displayed
inclusion of selectors. Thus the guard tables retain the minimal
suffixes needed for evaluation and comparison; they do not introduce
an unrelated normalization theory. The main text gives the
prerequisite/antimatroid presentation and its algorithms. Priority for
those connections and algorithms remains unestablished.

\paragraph{Antimatroids.}
To verify the description of $e_{\mathcal F}(B)$ in the main text, take a feasible subset of $B\cup\{j\}$ containing $j$ and an
accessible order of it. Its prefix before $j$ is feasible inside $B$; unioning with the basis makes $j$ an available extension. The
converse is immediate.

\begin{proof}[Proof of \cref{reset:thm:anti}]
The variable $x_j$ gives $\{\varnothing,\{j\}\}$.
The family in \eqref{reset:eq:family-product} is union-closed. To order a
nonempty member $A\cup C$ accessibly, first take an accessible order of
$C$ in $\mathcal G$, then one of $A$ in $\mathcal F$, skipping positions
already included. Every prefix is a union of a feasible set of
$\mathcal F$ with a nonempty feasible set of $\mathcal G$. This proves
accessibility, full support, and the unique initial element.

If $r$ is the root of $\mathcal G$, its basis obeys
\[
 \bas_{\mathcal F\del\mathcal G}(B)=
 \begin{cases}
 \varnothing,&r\notin B,\\
 \bas_{\mathcal F}(B)\cup\bas_{\mathcal G}(B),&r\in B.
 \end{cases}
\]
Applying the definition of $e$ gives exactly the selector recursion
in Lemma~\ref{reset:lem:selection}. Also
$A\subseteq e_{\mathcal F}(A)$ iff $A$ is feasible: membership of
$j\in A$ in this readiness set means $j\in\bas_{\mathcal F}(A)$.
This proves faithful representation for every term.

Conversely, an accessible order $r,j_2,\ldots,j_m$ of the full support
is a basic word. Its path term
\[
 x_{j_m}\del x_{j_{m-1}}\del\cdots\del x_r
\]
with \emph{left} parentheses has the feasible family of its prefixes.
Here the root $r$ is the last argument of the term, while it is the
first element of the accessible order. Combining two terms with the
same root by reset is pointwise union. On families it is union closure
of the two families: each already contains the common singleton root.
Every feasible set extends to a basic word, by accessibility and union
closure. Thus joining the path terms for all basic words gives exactly
the prescribed antimatroid. This is an application of the standard
basic-word representation, with reset supplying its composition.
\end{proof}

The parentheses in the path display are essential: for an order $r,a,b$, the term is $(x_b\del x_a)\del x_r$. It selects the initial
segment through the first resolved input. Repeated labels in more general terms are handled by \eqref{reset:eq:family-product};
replacing the occurrence tree by a graph on variable names can incorrectly combine paths from different occurrences.

\paragraph{Realizability.}

\begin{proof}[Proof of \cref{reset:thm:synthesis}]
For necessity, roots and self-independence give the first two
conditions. A minimal prerequisite $A$ enabling $j$ is feasible in the
represented antimatroid: readiness at $A$ has a feasible predecessor
inside $A$, and minimality makes that predecessor all of $A$.
Any accessible order of $A$ verifies constructibility.

For sufficiency, for each prerequisite $A\in\mathcal G_j$ choose its
greedy order $a_1,\ldots,a_m$ and make the path term with order
$a_1,\ldots,a_m,j$. Its last position is not repeated, by condition 2.
The first is the unique unconditionally enabled position $r$.
The endpoint has exactly prerequisite $A$. Every earlier position on
the path has a prerequisite contained in the preceding prefix, by
constructibility. Thus all the path's extra contributions are already
permitted by the candidate table. Join these paths by reset with their
common root. Every required prerequisite is present and no unwanted
contribution is introduced. There are $M$ paths of length at most
$n+1$.
\end{proof}

To see the equivalence of constructibility with the local form \eqref{reset:eq:complete}, note first that constructibility implies
completeness. Conversely, induct on the size of a prerequisite $A$. For each $k\in A$, choose $B\in\mathcal G_k$ inside $A$. Since
$k\notin B$, it is smaller than $A$ and constructible by induction. Append $k$ to its construction. Unions of constructible sets are
constructible: concatenate their orders, skipping repeated positions. Unioning these constructions reaches all of $A$. The implementation
checks the local condition and uses greedy orders for synthesis. An exponential antimatroid is therefore not needed for either equality or
synthesis.

The substitution formula \eqref{reset:eq:substitution} follows by substituting the inner output unions into the outer output union, and
substituting each inner root status into an outer guard.

\begin{proof}[Proof of \cref{reset:thm:minimal}]
A term selects all its occurring variables when every input is
ambiguous. Each of them is essential, as witnessed by assigning a proper
ambiguous set to all inputs except one assigned a larger set.
A nonprojection therefore has support of size at least two. Identify
its root variable with $y$ and every other occurring variable with $x$.
The resulting binary term has root $y$ and both arguments essential.
Its $x$ prerequisite cannot contain $x$, cannot be empty, and can only
use $y$. Thus it is $\{\{y\}\}$, while the $y$ prerequisite is empty.
This is exactly $x\del y$. Minimality follows.
\end{proof}

\subsection{Hardness}\label{app:proofs-hard}

\begin{proof}[Proof of \cref{gen:thm:complexity}]
For fixed $K$, a failed valuation is a polynomial certificate since the
algebra is finite and term evaluation is linear up to fixed constants.
For variable $K$, use the pair countervaluation theorem. A certificate
lists at most two role labels per input variable. Evaluate terms using
sparse sets: an owned node either collapses or unions its inputs, while
a permutation preserves cardinality. Each intermediate set has size at
most twice the number of variable leaves below it plus the number of
winner leaves. This is polynomial in tree-term size, independently of
the exponential size of the carrier and even when $K$ is binary encoded.
Evaluation and comparison are therefore polynomial. Validity is in coNP.

For hardness, \cref{thm:raw-pairs} below gives a linear reduction
using just two distinct owned operations. It uses neither constants,
role switches, nor derived union. Fixed-$K$ hardness also supplies
uniform hardness.
\end{proof}

For the remaining reductions, size is measured in the stated native signature; an expansion into owned-choice syntax need not preserve that
size. Use $A\del B=B$ when $B$ is a singleton, and $A\cup B$ otherwise. Write $\operatorname{rc}(a_1,\ldots,a_n;b)$ for the right comb
$a_1\del(a_2\del\cdots\del(a_n\del b))$, and $\operatorname{lc}(b;g_1,\ldots,g_m)$ for the left comb
$(\cdots((b\del g_1)\del g_2)\cdots)\del g_m$. Both are $b$ when their lists are empty.

\begin{proof}[Proof of \cref{thm:owner-free-hardness}]
Membership in coNP follows either from a full valuation at fixed $K$,
or from the pair theorem and polynomial native evaluation. The pair
witness exists because both operations are owned-choice definable;
its existence does not require a small expansion of the input term.

For hardness, take a CNF $\varphi$ on $n\geq1$ Boolean variables.
Use literal variables $u_i,\bar u_i$ and three additional input variables
$t,z,c$. Put
\begin{align*}
 W&=\operatorname{rc}(u_1,\bar u_1,\ldots,u_n,\bar u_n,t;z),\\
 b_i&=(c\del(u_i\cup t))\del(\bar u_i\cup t),\\
 R&=\operatorname{rc}(b_1,\ldots,b_n;W),\\
 g_C&=t\cup\bigcup_{\ell\in C}\ell,\qquad
 A=\operatorname{lc}(c;(g_C)_{C\in\varphi}).
\end{align*}
The equation is
\begin{equation}\label{eq:union-reduction}
 A\del R=R.
\end{equation}
All displayed symbols are input variables or the two native operations.
The tree size is linear in the CNF description and $n$.

If $z$ is resolved, $W$, $R$, and both sides of
\eqref{eq:union-reduction} are $z$. Otherwise $W$ is the union of
$z,t$, and all literals; $R$ is ambiguous and contains $W$.
Every contribution other than $c$ is therefore already in $R$.
Outside $W$, a marker from $c$ survives in $A$ exactly when every
clause guard $g_C$ is ambiguous. It survives in $R$ exactly when
some pair $u_i\cup t,\bar u_i\cup t$ is simultaneously ambiguous.

Suppose the equation fails. Every clause guard is ambiguous, and no
pair is simultaneously ambiguous. Since $n\geq1$, $t$ must be a
singleton: an ambiguous $t$ would make every pair ambiguous.
For nonempty $X$, $X\cup t$ is resolved exactly when $X=t$.
Assign variable $i$ true iff $u_i\ne t$. If a positive literal differs
from $t$, it is true under this assignment. If a negative literal
$\bar u_i$ differs from $t$, consistency forces $u_i=t$, so that
negative literal is true. Every ambiguous clause guard has such a
literal, hence $\varphi$ is satisfied.

Conversely choose three distinct roles $b,d,e$ and a satisfying
assignment. Set $t=\{b\}$, each false literal to $\{b\}$, each true
literal to $\{d\}$, $z=\{b,d\}$, and $c=\{e\}$.
All clause guards are ambiguous, while every literal pair has a
resolved guard. Thus $e$ survives on the left and is absent on the
right. This proves the reduction. CNFs with no variables are decided
trivially and mapped to a fixed valid or invalid equation.
\end{proof}

\begin{proof}[Proof of \cref{thm:H-hardness}]
The same finite-valuation argument gives membership in coNP.
For a CNF $\varphi$, use one input $y_i$ per Boolean variable, together
with $z,c$. Let $W=\operatorname{rc}(y_1,\ldots,y_n;z)$.
For a clause $C=(\ell_1,\ldots,\ell_m)$ define
\[
 b_C=H_{s_1}(z,y_{i_1},H_{s_2}(z,y_{i_2},\ldots
                  H_{s_m}(z,y_{i_m},c)\ldots)),
\]
where $s_j=a$ for a positive literal and $s_j=b$ for a negative one.
An empty clause gives $b_C=c$. Put
$R=\operatorname{rc}((b_C)_{C\in\varphi};W)$. The equation is
\begin{equation}\label{eq:H-reduction}
 c\del R=R.
\end{equation}
Expand each reset by $x\del y=H_a(y,x,x)$. In each right comb, the
continuation $y$ occurs once and only the small left operand is
repeated. The clause terms also use their continuation once. The
expanded tree therefore has polynomial, in fact linear, size in
$n$ and the total literal count.

When $z$ is resolved, both sides are $z$. Otherwise $W$ and $R$ are
ambiguous, and $W$ contains $z$ and every $y_i$. Outside $W$ the
exact propagation law is
\[
 H_s(z,y,u)\text{ contains a marker from }u
 \quad\Longleftrightarrow\quad
 u\text{ contains it and }y\ne\{s\}.
\]
This also holds when the marker is $s$: then $y\ne\{s\}$ follows
already from $y\subseteq W$. It follows by induction that a marker
from $c$ survives through $b_C$ iff all its literal tests
$y_{i_j}\ne\{s_j\}$ hold.

If $\varphi$ is unsatisfiable, read variable $i$ as true iff
$y_i=\{a\}$. Some clause is false. Its positive literals have
$y_i\ne\{a\}$, and its negative literals have $y_i=\{a\}\ne\{b\}$.
Consequently that branch retains every marker outside $W$ from $c$.
Since $W\subseteq R$, this gives $c\subseteq R$ and proves
\eqref{eq:H-reduction} on every valuation.

For a satisfying assignment, set $y_i=\{a\}$ for true and
$y_i=\{b\}$ for false, $z=\{a,b\}$, and $c=\{d\}$ for a third
role $d$. Every clause has a literal with $y_i=\{s\}$, removing
$d$ in its branch. Thus $d\notin R$, while $d\in c\del R$.
\end{proof}

\begin{proof}[Proof of \cref{thm:colouring-theories}]
For $n\geq3$ use variables $x_1,\ldots,x_n,z,c$. Define
\begin{align*}
 W_n&=\operatorname{rc}(x_1,\ldots,x_n;z),\\
 R_n&=\operatorname{rc}(c\del x_1,\ldots,c\del x_n;W_n),\\
 A_n&=\operatorname{lc}(c;(x_u\cup x_v)_{1\leq u<v\leq n}).
\end{align*}
Take the equation $A_n\del R_n=R_n$.
If $z$ is resolved, it holds immediately. Otherwise $W_n$ is ambiguous
and contains all the $x_i$. If any $x_i$ is ambiguous, $R_n$ also
contains $c$, so the equation holds. If all $x_i$ are singletons and
two agree, their edge guard is resolved and erases $c$ from $A_n$;
every remaining contribution lies in $W_n\subseteq R_n$.
The only remaining case has $n$ distinct singleton inputs. At $K=n$
they exhaust the carrier's roles, so $W_n=I$ and the equation holds.

At $K=n+1$ set $x_i=\{i\}$, $z=\{1,2\}$, and $c=\{n+1\}$.
Every edge guard is ambiguous; $A_n$ retains the fresh winner $n+1$,
while $R_n=W_n=\{1,\ldots,n\}$ omits it. The equation fails.
Finally, the inclusion of role sets induces a subalgebra embedding
$\mathcal B_K\hookrightarrow\mathcal B_{K+1}$: union and singleton
status are preserved. Equations valid in the latter hold in the
former, and the displayed equation makes this inclusion strict.
\end{proof}

\begin{proof}[Proof of \cref{thm:raw-pairs}]
A single owned operation and union are semilattices: equality is equality
of variable supports. Reset has the polynomial algorithm of
\cref{reset:thm:three}. The empty signature permits only variables.
Nonvalidity has a polynomial valuation certificate at fixed $K$.
For hardness we reduce CNF unsatisfiability, using the standard
SAT completeness theorem \cite{Cook1971}.
It remains to cover four types of distinct pairs. Reset with union is
\cref{thm:owner-free-hardness}.

For two owners $a\ne b$, write $\odot=\star_b$. Given a CNF with at
least one clause, all clauses nonempty, and $n\geq1$, use rails
$p_i,n_i$ and marker variable $m$. Put
\begin{align*}
 W&=p_1\odot n_1\odot\cdots\odot p_n\odot n_n,\\
 b_i&=m\star_a p_i\star_a n_i,\qquad
 R=b_1\odot\cdots\odot b_n\odot W,\\
 u_C&=\mathop{\bigodot}_{\ell\in C}\ell,\qquad
 A=m\star_a u_{C_1}\star_a\cdots\star_a u_{C_q}.
\end{align*}
The equation is $A\odot R=R$. Associativity makes the displayed folds
unambiguous. Empty clauses and empty conjunctions are handled by fixed
valid and invalid equations respectively.

For a satisfying Boolean assignment, choose $c\notin\{a,b\}$, assign
true rails $\{c\}$, false rails $\{a\}$, and $m=\{b\}$.
Every $b_i$ is $\{a\}$ and $R=\{a,c\}$, whereas each clause value
contains $c$ and $A$ contains $b$. The equation fails.

Conversely suppose the CNF is unsatisfiable. If a rail is $\{b\}$,
$W=R=\{b\}$ and absorption proves the equation. Otherwise every
clause fold is a union and differs from $\{b\}$. Each $b_i$ also
differs from $\{b\}$: an $a$-fold equals that different singleton only
when all its inputs do. Thus $R$ is a union containing $W$, and
$A\ne\{b\}$. It suffices to prove $A\subseteq R$.
If $A=\{a\}$, either $m=\{a\}$, making every $b_i=\{a\}$, or
some $u_C=\{a\}$, already contained in $W$. Otherwise $m$ and all
$u_C$ differ from $\{a\}$ and $A=m\cup\bigcup_C u_C$.
If a pair of rails both differs from $\{a\}$, its $b_i$ retains $m$.
If no such pair exists, assign each variable so that its false literal
has rail $\{a\}$. Unsatisfiability supplies an all-false clause,
whose fold is $\{a\}$, a contradiction. Hence $m\subseteq R$,
and the equation holds.

For $\{\star_a,\cup\}$ use exactly the same construction with
$\odot=\cup$. The $\{b\}$ cases are unnecessary; the containment
argument and the same three-role countervaluation apply.

For $\{\del,\star_a\}$ use \cref{thm:owner-free-hardness}'s construction,
replacing every explicit union by $\star_a$. The reset comb $W$ still
contains all rails and $t$ whenever $z$ is ambiguous. Every new guard
is contained in that baseline. An outside marker fails the equation
only if all clause guards are ambiguous while no rail pair has both
its guards ambiguous. For any nonempty clause,
\[
 \left|t\star_a\mathop{\star_a}_{\ell\in C}\ell\right|>1
 \quad\Longrightarrow\quad
 \exists\ell\in C:\ |\ell\star_a t|>1.
\]
To verify this implication, an ambiguous fold has no input equal to
$\{a\}$. If $t$ is ambiguous, any clause literal gives the witness.
If $t=\{r\}\ne\{a\}$, the fold is a union and some literal differs
from $\{r\}$, giving the witness. The unsatisfiable-rails argument now
forces a conflicting pair, contrary to failure.
For a satisfying assignment choose distinct $a,b,d$ and set
$t=\{b\}$, false rails $\{b\}$, true rails $\{d\}$,
$z=\{b,d\}$, and $c=\{a\}$. No guard operand is $\{a\}$,
so every replaced union agrees with its old value and the marker
separates the equation as before. This proves the final pair case.
All constructions use each clause and rail only a bounded number of
times. Hardness passes to every larger signature.
\end{proof}

\subsection{One owner}\label{app:proofs-owner}

\begin{proof}[Proof of \cref{owner:lem:root}]
Reset returns its right operand when that operand is resolved, and
otherwise returns an ambiguous union containing it. Both status claims
follow by induction, since the right operand in the definition of
$H_r(X,Y,Z)$ is $X$.
\end{proof}

\begin{proof}[Proof of \cref{owner:thm:guards}]
Sufficiency follows from \eqref{owner:eq:guards} and
\eqref{owner:eq:cover}. For necessity, choose $b\notin\{k,a\}$ and
an uncovered guard $g=(j,A,B)$. Assign
\[
 x_i=\begin{cases}
 \{k,a\},&i=j\text{ and }i\in A,\\
 \{k\},&i=j\text{ and }i\notin A,\\
 \{a,b\},&i\ne j\text{ and }i\in A,\\
 \{b\},&i\ne j,\ i\notin A,\ i\in B,\\
 \{a\},&\text{otherwise}.
 \end{cases}
\]
Exactly the indices in $A$ are ambiguous; exactly those in
$A\cup B\cup\{j\}$ satisfy $N$; only input $j$ contains $k$.
Thus $g$ holds, and a right guard holds exactly when it covers $g$.
No right guard does. This is the required countervaluation.
\end{proof}

\paragraph{Canonical tables.}
The criterion also gives canonical tables: for each source, remove its index and the $M$ indices from its $N$ tests, then retain only the
minimal guards under \eqref{owner:eq:cover}. No minimal guard has its source in its $M$ tests. Indeed, a guard $(j,A,B)$ with $j\in A$ has a
covering guard $(j,A',B')$ with $j\notin A'$, $A'\subseteq A$, and $B'\subseteq B$. Induct on the term: if its root is $j$, the root
occurrence supplies $(j,\varnothing,\varnothing)$; otherwise apply the induction hypothesis to the contributing child and reattach the node's
tests. Any newly attached $M$ test is the node's different root. Minimal antichains are unique because \cref{owner:thm:guards} characterizes
implication exactly. Guard compilation and the three tests in \eqref{owner:eq:cover} check the short certificate of an affirmative comparison
directly.

\begin{proof}[Proof of \cref{owner:thm:four}]
First retain the five flags used to evaluate the recursion:
\[
 (m,e,p)=(|S|>1,\ S=\{a\},\ k\in S).
\]
The five possible flag states form a homomorphic image of the winner-set
algebra. Identify the two states with $p=1$: observed singletons and
observed ambiguous sets. This is a congruence. In a root position both
already give observed victory; in every other position the rules use
only $e$ and $p$, which agree. The four resulting states have exactly
the displayed operations on $\mathbf C$.
An identity of $\mathbf C$ therefore preserves the outsider's $p$ flag
on all winner-set valuations. By \cref{owner:thm:one-table} it is a full
identity. Conversely, $\mathbf C$ is a subalgebra, so it satisfies every
winner-set identity. This proves completeness and independence of $K$.

For minimality, any smaller generator has a nonassociative reset reduct.
By Freese et al.'s Theorem~3.2 \cite{FreeseEtAl2002}, such a reduct
contains $O_{2+1}$; hence a three-element candidate is precisely that
algebra. Name its elements $Z,U,T$, with
\[
 u\del Z=Z,\quad u\del T=T,\quad
 Z\del U=U\del U=U,\quad T\del U=T.
\]
Write $h$ for the candidate ternary operation. With $s=h(x,y,c)$,
the valid identities $s\del x=s=x\del s$ force
$h(Z,y,c)=Z$, $h(T,y,c)=T$, and $h(U,y,c)\in\{U,T\}$.
The following identities are also valid in $\mathbf C$:
\[
 h(x,y,y)=y\del x,\qquad h(y,y,c)=c\del y,\qquad
 h(x,y,x)=y\del x.
\]
Put $t=(c\del y)\del x$. The equation $s\del t=t$ is false in
$\mathbf C$, at $(x,y,c)=(N,B,P)$. On the three-element chain,
the root bounds and these identities allow it to fail only when
$h(U,Z,T)=T$.
The equation $(c\del x)\del s=s$ is also false in $\mathbf C$, at
$(N,A,P)$. In a three-element candidate its failure requires
$h(U,y,T)=U$ for some $y$. The binary-minor identities give value $T$
when $y=U,T$, while the first required failure gives $T$ when $y=Z$.
This is impossible. Smaller reducts are associative, so they cannot
generate the theory either.
\end{proof}

\begin{proof}[Proof of \cref{owner:cor:complexity}]
\Cref{owner:thm:one-table,owner:thm:guards} give the algorithms; unused
owner indices may be added to apply them. Two distinct $H$ operations
give the full-equation lower bound by \cref{thm:H-hardness}.
For the observational lower bound choose two owners $a,b\ne k$ in that
same reduction and use $c=\{k\}$ as its fresh-marker witness. In an
unsatisfiable instance the equation is valid for all roles; in a
satisfiable instance this valuation separates membership of $k$.
The reduction is to the implication $c\del R\subseteq_k R$, where
$\subseteq_k$ means implication of $k$ membership, as well as to its
equivalence. Nonvalidity has a polynomial singleton-or-pair certificate,
so both problems belong to coNP.
\end{proof}

\subsection{Sealed rounds}\label{app:proofs-sealed}

\begin{proof}[Proof of \cref{se:thm:round}]
The standard finite-horizon one-stage-deviation argument gives the
mixed SPE recursion. We prove its equivalence with the serial SE model.
After full publication, each $G_a$ is a proper public subgame. Restricting
an assessment to it gives a sequential equilibrium and a distribution
$p_a$. Any pure change of the current action, followed by the prescribed
continuation, is a legal deviation. The current independent action laws
therefore form a Nash equilibrium for the displayed payoffs.

Conversely choose one continuation equilibrium at every tuple and a
stage Nash profile. Backward induction rules out deviations that also
change later actions: each selected continuation is already optimal
against its opponents, and a current action change cannot improve its
expected continuation payoff. To verify consistency rather than just
subgame optimality, perturb every action law toward its uniform law
with one common positive parameter. At a public round history its reach
factor cancels from Bayes' rule. At a player's commitment, the belief
over preceding hidden commitments is the product of those roles'
perturbed laws. Its limit is the product of the prescribed laws, also
at off-path public histories. Revelation steps offer no choices. This
constructs one consistent assessment for the whole finite game.
The same argument works for every serial order of private commitments.

Nonemptiness follows by choosing continuation distributions and applying
finite mixed Nash existence \cite{Nash1951}. The variables $p_a$ and
$\sigma_i$ range over compact sets. Nash optimality is a finite family
of weak polynomial inequalities, so the feasible set is closed and
its continuous image is compact. The simplex, these inequalities, and
the output equation are rational polynomial conditions. Projection
preserves semialgebraicity, giving the induction. Inclusion of
continuation sets preserves every old feasible choice. Role renaming
only permutes the inequalities and probability coordinates.
\end{proof}

\begin{proof}[Proof of \cref{se:cor:public}]
Pure SPE profiles are sequential equilibria in this perfect-information
language. For the other inclusion induct on trees. If a child has pure
winner set $\{r\}$, its mixed continuation distributions are all
$\delta_r$ by induction. The parent owner $r$ can secure utility one,
so every parent equilibrium also has distribution $\delta_r$.
Otherwise the pure recurrence already includes every possible child
winner; mixing can introduce no others. Nonemptiness of each child's
equilibrium set is used throughout.
\end{proof}

\begin{proof}[Proof of \cref{se:prop:parallel}]
For $x+y>0$, both players assign probability $y/(x+y)$ to their
action $0$. Against this law, either pure action of the opponent gives
$xy/(x+y)$. These two security guarantees force that value in every
Nash equilibrium. If one diagonal entry is zero the same value follows;
if both are zero every strategy profile has value zero.
For positive arguments, reciprocal values turn parallel sum into
addition. This proves associativity and commutativity; the zero cases
are immediate. Monotonicity follows from the displayed formula, and
a monotone unary map on a chain preserves binary minima and maxima.
\end{proof}

\begin{proof}[Proof of \cref{se:thm:rational}]
The operations preserve rationality. For the converse induct on the
denominator of $p/q$, with $0\leq p\leq q$ and $q>0$. Endpoints are
leaves. If $0<p<q$ and $2p\leq q$, write
\[
 p/q=U\bigl(p/(q-p)\bigr).
\]
Otherwise write
\[
 p/q=C\bigl(U((q-p)/p)\bigr).
\]
Each recursive denominator is smaller. Reduction of fractions is not
needed. This gives an explicit synthesis procedure, without claiming
polynomial size in the binary encoding of $p,q$.
\end{proof}

\begin{proof}[Proof of \cref{se:cor:decidable}]
All primitives are continuous. Parallel sum is continuous also at
$(0,0)$ because $0\leq x\parallel y\leq\min(x,y)$.
Rational valuations are dense. For the decision claim, the graph of
parallel sum is given on the unit square by $z(x+y)=xy$, with $z=0$
when $x+y=0$; the other graphs are polynomial or piecewise polynomial.
Term evaluation and a failed equality are therefore expressible in the
first-order theory of real closed fields. Its decision theorem applies.
\end{proof}
